\documentclass[11pt]{article}
\usepackage[margin=1in]{geometry}
\usepackage{amsmath}
\usepackage{amssymb}
\usepackage{graphicx}
\usepackage{booktabs}
\usepackage{xcolor}
\graphicspath{{figures/}%
  {../output/figures_sllg/track_width/box350x500_racetrack_D0p545/}}

\newcommand{\Q}{\lvert Q\rvert}

\title{Data-driven stability classification of the racetrack
skyrmion campaign\\[2pt]
\large box $350\times500$, racetrack BC, $D=0.545$~mJ/m$^2$}
\author{Rui Barreira}
\date{2026-07-12}

\begin{document}
\maketitle

\section{Motivation}
The first stability map was produced by a hardcoded two-threshold
lookup on $(T,J)$: labyrinth for $T\ge130$~K, elongated stripe for
$J\ge3\times10^{11}$~A/m$^2$, else compact skyrmion. Those thresholds
were read by eye off the snapshots of a \emph{different} run (the
periodic $350\times500$ \texttt{pbcx\_pbcy} box) and then frozen; they
never consulted the actual configurations of the racetrack $D=0.545$
run, and carried no guard against the periodic-boundary artifact below.
This note replaces the lookup with a classifier computed from the raw
simulation output, and documents the method and results.

\paragraph{The periodicity artifact.}
The track is periodic along $x$ (motion axis, length $L_x=n_x a=700$~nm)
and free along $y$. Once the reversed domain extends along $x$ to about
half the track, $D_x\gtrsim0.5\,L_x$, it can connect to its own
periodic-$x$ image. The largest-connected-component (LCC) analysis then
describes a wrapped, image-bridged object, and the single-skyrmion
picture is invalid there. Crucially the relevant length is the
\emph{$x$-extent} $D_x$, \emph{not} the ellipse major axis $D_1$: a
skyrmion stretched along free-$y$ has large $D_1$ but small $D_x$ and
does not bridge its $x$-image. We measure $D_x$ directly from the field
and flag the affected cells explicitly rather than trust them.

\section{Data reuse --- no new simulation}
Everything needed already exists in the completed campaign output; no
resimulation was required.
\begin{itemize}
  \item \textbf{Topological charge $\Q$}: stored per frame
  (\texttt{Q\_top}) in the 36 \texttt{anim\_*.npz} full-field dumps
  (the \texttt{ens0} representative of each $(T,J)$ cell).
  \item \textbf{Domain structure and $D_x$}: computed from the final
  drive-frame $m_z(x,y)$ field in the same dumps (the only place full
  fields were saved).
  \item \textbf{$D_1/D_2$ elongation ratio}: the ensemble-mean ellipse
  axes in \texttt{aggregate\_350x500\_racetrack\_D0p545.npz}.
\end{itemize}

\section{Methodology}
Each $(T,J)$ cell is classified from the final drive-phase
configuration of its \texttt{ens0} realization, combining three
physical signals.

\paragraph{1. Topology.}
$\Q$ is the integrated topological charge: $\Q\approx1$ marks one
skyrmion; $\Q\approx0$ a collapse to the ferromagnet or a chargeless
stripe; $\Q\gtrsim1.6$ more than one skyrmion.

\paragraph{2. Domain structure (de-speckled).}
The reversed domain is the mask $m_z<-0.5$ --- a hard threshold, so
thermal fluctuations near $m_z=0$ are not miscounted as domains.
Connected components are labelled with periodic-$x$ / free-$y$
connectivity (a union-find merge across the $x$ seam), and components
below $0.2\%$ of the box are dropped as speckle. From the survivors we
take the count $N_c$, the reversed-area fraction, whether any component
wraps the $x$ seam, and the periodic $x$-extent $D_x$ of the largest
component (span $=n_x-$ largest cyclic gap of empty columns).

\paragraph{3. Periodic gate.}
$D_x/L_x\ge0.5$ means the domain can bridge its periodic image (loss of
periodicity).

\paragraph{Decision order.}
\begin{enumerate}
  \item No surviving domain (or a sub-floor domain with $\Q<0.5$)
  $\rightarrow$ \textbf{A} (annihilated, ferromagnetic).
  \item $N_c\ge2$, or reversed fraction $\ge0.35$, or $\Q\ge1.6$
  $\rightarrow$ \textbf{L} (labyrinth / multi-domain).
  \item Single domain with $\Q<0.5$: \textbf{E} if it spans the seam,
  else \textbf{L}.
  \item One skyrmion ($\Q\approx1$), single domain: \textbf{E} if it
  wraps, $D_x/L_x\ge0.5$, or $D_1/D_2\ge1.7$; else \textbf{S}
  (compact skyrmion).
\end{enumerate}
Independently, a second index \textbf{p} is appended wherever
$D_x/L_x\ge0.5$ (loss of periodicity): e.g. \texttt{Ep} is an elongated
state whose $x$-size is no longer trustworthy, \texttt{Lp} a labyrinth
in the image-bridged regime.

\begin{table}[h]
\centering
\caption{Classifier thresholds (physical constants, local to the
routine).}
\begin{tabular}{lll}
\toprule
quantity & value & meaning\\
\midrule
core threshold & $m_z<-0.5$ & reversed domain vs thermal noise\\
speckle floor & $0.2\%$ of box & minimum component area\\
$\Q$ skyrmion & $0.5$ & topological skyrmion present\\
$\Q$ multi & $1.6$ & more than one skyrmion\\
reversed fill & $0.35$ & space-filling $\rightarrow$ labyrinth\\
$D_x/L_x$ gate & $0.5$ & periodic-image bridging (index p)\\
$D_1/D_2$ & $1.7$ & elongated\\
\bottomrule
\end{tabular}
\end{table}

\section{Results}
\begin{figure}[h]
\centering
\includegraphics[width=0.82\linewidth]{stability_map.png}
\caption{Stability regimes over the $(T,J)$ grid. Colour is the
physical class (S/E/L/A); the appended \textbf{p} marks loss of
periodicity ($D_x\ge0.5\,L_x$).}
\label{fig:stab}
\end{figure}

\begin{table}[h]
\centering
\caption{Classification grid (rows $T$ in K, columns $J$ in
$10^{11}$~A/m$^2$). Code $=$ physical class $+$ optional \textbf{p}
(loss of periodicity, $D_x\ge0.5\,L_x$).}
\begin{tabular}{r|cccccc}
\toprule
$T\,\backslash\,J$ & 0.5 & 1 & 2 & 3 & 4 & 5\\
\midrule
10  & S  & S  & E  & E  & E  & Lp\\
50  & S  & E  & E  & E  & Lp & Lp\\
100 & Lp & Lp & Ep & Ep & Ep & Ep\\
130 & Ep & Ep & Lp & Ep & Lp & Lp\\
160 & Ep & Lp & Lp & Lp & Lp & Lp\\
200 & L  & L  & L  & Lp & L  & Lp\\
\bottomrule
\end{tabular}
\end{table}

\paragraph{Findings.}
\begin{itemize}
  \item A genuinely compact skyrmion (\textbf{S}) survives only in the
  cold, low-drive corner: $T=10$~K at $J\le1\times10^{11}$ and $T=50$~K
  at $J=5\times10^{10}$. There $\Q=1.00$, $D_1/D_2\approx1.0$--$1.1$,
  $D_x/L_x\approx0.2$--$0.5$.
  \item With increasing current the low-$T$ skyrmion stretches
  (\textbf{E}: $D_1/D_2$ rising to $3.2$ at $T=10$), then collapses to a
  labyrinth at the highest current (\texttt{Lp}; $\Q=0.17$ at $T=10$,
  $J=5\times10^{11}$).
  \item \textbf{Loss of periodicity affects $23$ of $36$ cells}, but the
  field-measured $D_x$ gate is far more discriminating than the ellipse
  $D_1$: the round $T=50$, $J=5\times10^{10}$ skyrmion
  ($D_1/D_2=1.11$, $D_x/L_x=0.47$) is correctly \textbf{S}, whereas a
  $D_1$-based gate flagged it (its $D_1$ was thermally inflated). At
  warm/high-drive cells $D_x$ reaches or exceeds $L_x$.
  \item Warm cells ($T\ge130$) are labyrinthine at most currents; at
  $T=200$ the texture fragments into several small domains, so the
  \emph{largest} component's $D_x$ is small and most cells there are
  plain \textbf{L}.
\end{itemize}

\paragraph{Consequence.}
The reliable, intrinsic region of this run is the small cold/low-drive
corner. To characterise the elongated and driven states without the
periodic-image artifact, the track must be extended along the motion
axis (larger $L_x$ at fixed width), the subject of a follow-up campaign
($256$, $512$, $1024$ cells along $x$, free $y$).

\end{document}
